\documentclass{article} % ICLR 2027 submission format (variant of NeurIPS)
\usepackage{iclr2027_conference,times}
\input{math_commands.tex}
\usepackage{hyperref}
\usepackage{url}
\usepackage{booktabs}
\usepackage{graphicx}
\usepackage{amsmath}

% ---------------------------------------------------------------------------
% ICLR 2027 rules (https://iclr.cc/Conferences/2027/AuthorGuidelines):
%   * abstract deadline Sep 18 2026 AoE, paper deadline Sep 25 2026 AoE
%   * main text <= 9 pages at submission (10 at camera-ready); references and
%     appendices unlimited
%   * double blind: no author names, no identifying URLs, cite own work in
%     third person; leave \iclrfinalcopy commented out for submission
%   * AI-use statement required; Reproducibility Statement strongly encouraged
% ---------------------------------------------------------------------------

\title{Rote: Managing the Context Window as a Resource \\ for Browser Agents that Learn Across Runs}

% Anonymous for submission. Uncomment \iclrfinalcopy for camera-ready only.
\author{Anonymous authors \\ Paper under double-blind review}

%\iclrfinalcopy

\begin{document}
\maketitle

\begin{abstract}
% ~150--200 words. Claim, mechanism, evidence, caveat.
Browser agents built on large language models spend most of their budget re-reading pages they have already seen and re-deriving actions they have already taken. Every agent harness \emph{has} memory; almost none \emph{manages} it. We present Rote, a memory manager that treats the context window as a managed resource---with an explicit budget, an eviction policy, a layout contract that keeps a cache-stable prefix, and a trust gate on everything re-admitted---across three tiers: working memory within a run, episodic memory across runs of the same task, and semantic memory across tasks on the same site. Rote distills recorded trajectories into value-free, contract-gated playbooks; matches new tasks against them with a fail-closed matcher; and replays with zero model calls until a versioned \emph{action contract} detects that a same-looking control changed its behavior, at which point it falls back to the plain agent before any dispatch. On a frozen enterprise-style benchmark, Rote reduces logical input tokens by 83.1\% (95\% CI 82.1--83.9\%) and billed cost by 68.2\% at exact-oracle parity with a strong open-source baseline, and its offline next-action predictor reaches 99.4\% kind+target accuracy over 1{,}520 warm steps. All mechanisms fail closed: 0/90 silent failures were observed under selector drift, destructive decoys, and ambiguity. \textbf{[TODO: replace with numbers from the billed learning-curve campaign once run.]}
\end{abstract}

\section{Introduction}
% 1 page. Problem: agents re-pay for what they already know; memory is unmanaged.
% Contribution list (keep to 4):
%  C1 Context window as managed resource: budget, eviction, layout contract, trust gate (docs/02).
%  C2 Learn-across-runs loop with hard safety: distill -> match -> replay, value-free, append-only,
%     fingerprint hard gate, action-contract gate (#143 / T35, T36, T40).
%  C3 Predictor as a shadow system with monotone confidence; kill gate + calibration (T38, T39).
%  C4 Evidence: frozen benchmark B1--B6, exact external oracles, invariant suite; billed campaign.
\textbf{[TODO]}

\section{Related work}
% Web/browser agents (WebArena, Mind2Web, Browser Use, Stagehand, Skyvern, Magnitude -- we qualified
% three of these as baselines, T22/T23/T24/T27); agent memory (MemGPT-style paging, reflexion/
% experience replay, skill libraries e.g. Voyager); prompt caching / KV-cache reuse; program synthesis
% from demonstrations (macro recording, PbD). Position: prior work adds memory; Rote manages it with
% invariants. Cite own arXiv (if any) in third person.
\textbf{[TODO]}

\section{Rote: the memory spine}
\label{sec:design}

Rote is organized around a single reading of the agent loop: the context window is
memory, and memory needs a manager --- a budget, an eviction policy, a write path, and a
trust gate on everything re-admitted. The operating-system analogy is exact and we treat
it as a design discipline rather than a metaphor: the context window is RAM, observations
are pages, dropping them is eviction, diffing is delta encoding, the provider prompt
cache is a cache level, compaction is garbage collection, a distilled playbook is a
cached compiled program, and site memory is the persistent store. Each of these has a
manager in an operating system; none has one in a shipping browser agent, which appends
to a transcript and hopes. Rote supplies the manager at three timescales: \emph{working}
memory within a run (\S\ref{sec:tier0}), \emph{episodic} memory across runs of one task
(\S\ref{sec:tier1}), and \emph{semantic} memory across tasks on one site
(\S\ref{sec:tier2}). The trust gates (\S\ref{sec:gates}) are not a fourth tier but the
precondition for all three: memory that might be wrong is worse than no memory, because
reuse without verification is a machine for repeating a mistake at volume.

A second commitment follows from the first: control flow should be deterministic. A
successful run tangles together \emph{what to do} (the procedure) and \emph{what to
decide} (the content). Rote untangles them --- the procedure becomes a deterministic,
replayable artifact, and the model is invoked only where genuine judgment lives: binding
parameters, filling slots, and repairing broken steps.

\subsection{Tier 0: working memory as a budgeted resource}
\label{sec:tier0}

The dominant cost in a transcript-style agent is quadratic. With per-step action records
of $\sim$35--40 tokens, step $n$ costs $c + 40n + \mathit{obs}$, so an $n$-step run
carries a $40\cdot n(n+1)/2$ term; if every observation also stays resident at 5--40K
tokens each, step 20 is a 100--800K-token prompt. Rote's context assembler enforces a
layout contract with four segments --- a constant instruction prefix, a constant page
header, a linearly growing action history, and the \emph{current} observation only ---
and applies four levers against the curve. (1)~\textbf{Observation eviction}: prior
observations are dropped; the model recalls what it \emph{did}, not what it \emph{saw}.
This is a real trade: a task whose answer lives in an evicted observation cannot be
completed, so after first omission the context marks the recall boundary and instructs
the planner to return a typed \texttt{recall\_unavailable} failure instead of guessing;
independent verification separately rejects fabricated answers. The task still fails ---
the protection is that it fails honestly. (2)~\textbf{Observation diffing}: the current
observation is served as a delta against a grounded base; across 849 real-page
certification diffs the median delta is 24 characters, a 99.6\% median reduction.
(3)~\textbf{Cache-stable layout}: the prompt is ordered $[\text{stable
prefix}][\text{append-only history}]$ and routed by a deterministic immutable-prefix
cache key, so the surviving prefix is billed at the provider's cached rate; measured on a
25-step cell this cuts billed cost by 20.5\%. (4)~\textbf{Scheduled compaction}: past a
budget threshold the action history is folded into a versioned summary with the recent
suffix kept verbatim, holding visible history $O(1)$ in steps; a 60-transition
single-page-application session sustains this with at most 31 visible actions and three
compaction boundaries.

\subsection{Tier 1: episodic memory --- record, distill, replay}
\label{sec:tier1}

Every run is recorded to an append-only, crash-safe event log (one fsync per event, large
results spilled to content-addressed blobs). The \emph{distiller} converts one successful
recorded trajectory into a playbook: a parameterized, value-free procedure in which
concrete user inputs are lifted to named parameters, and every step carries the target's
structural identity, its \emph{action contract} (\S\ref{sec:gates}), and a learned
verification derived from the recorded evidence class of the original success. Values
never survive distillation --- credentials and form contents are replaced by parameter
references at distill time, so a playbook can be shared or inspected without leaking a
recorded secret. Replay executes the playbook deterministically with zero model calls,
asserting each step's expectations before and after dispatch; on our benchmark tasks a
recorded run distills and replays end-to-end with zero edits and zero tokens. A failed
assertion never degrades into a guess: the run stops, the failure is classified, and the
system falls back to the plain agent with the classification attached.

\subsection{Tier 2: semantic memory --- site memory and the brief}
\label{sec:tier2}

Across tasks on one site, Rote accumulates \emph{value-free} records keyed by a stable
page digest: selector maps, form field semantics, page-transition edges, settle-time
priors, and a closed vocabulary of site quirks. Value-freedom is enforced structurally
--- record schemas have no field that can hold page content, and an invariant test
derives site memory from a live run and asserts no user value appears anywhere in the
store. At run time the relevant records are rendered into a \emph{site brief}: a
hard-capped, deterministic block placed in the cache-stable prefix and framed as
advisory hints, never as instructions the planner must obey. Every run records the
brief's size and how many of its hinted identities were actually used, so the brief's
utility is measured rather than assumed.

\subsection{Trust gates}
\label{sec:gates}

Four gates stand between stored memory and a dispatched action, ordered from cheapest to
most specific. (1)~\textbf{Environment fingerprint} (hard gate): a structural fingerprint
of the environment must match exactly before any fuzzy matching is attempted; a playbook
learned on a staging host can never fire on production. (2)~\textbf{Matcher} (fail
closed): a new task selects a playbook only if intent similarity clears a threshold
($\tau = 0.8$) with a minimum margin over the runner-up (0.05), every required parameter
binds, and every content-bearing token of the task's intent is covered --- the coverage
requirement was added after a near-miss in which a destructive lookalike task scored
0.82 against a benign playbook. On any ambiguity the matcher declines and the plain
agent runs. (3)~\textbf{Action contract}: each step stores a versioned, value-free
contract of the control it dispatches --- element role, input modality, and the
observed effect class of activation. Before dispatch, the live control's derived
contract must equal the stored one; a same-looking button whose behavior changed (a
destination change, a method change, a purge action behind a familiar label) stops the
replay with a typed mismatch and zero dispatches, while cosmetic drift, selector
renames, and wrapper insertion replay normally. (4)~\textbf{Outcome evidence}: task
success is asserted only against authoritative, freshness-bound evidence tied to the
task subject --- a harness's own conclusion, stale evidence, or evidence about a
different subject all fail typed.

\subsection{Predictor and model routing}
\label{sec:routing}

With memory managed, two further systems spend the remaining model budget deliberately.
The \emph{predictor} is an ensemble of trace matching against prior runs and a learned
transition model, emitting a next-action prediction with a monotone confidence score. It
runs in \emph{shadow mode}: predictions are recorded and scored on every run but never
dispatched, because offline confidence is not yet calibrated on live pages
(\S\ref{sec:results-predictor}). The \emph{router} sends confident steps to a routine
(smaller) planner model and escalates to the frontier model on planner output errors or
target repair, preserving the never-worse-than-baseline invariant; escalation spend is
accounted separately so routing economics are measurable. Every model call in the system
carries a source tag through a single client wrapper (\S\ref{sec:invariants}), so all
reported token numbers decompose exactly by subsystem.

\section{Invariants}
\label{sec:invariants}

Five invariants bind every mechanism in \S\ref{sec:design}. They are encoded in a
dedicated invariant test suite that must touch every executor exit path, and they are
not negotiable under schedule pressure; several results in \S\ref{sec:results} are best
read as these invariants holding under adversarial pressure.

\begin{table}[h]
\caption{Project invariants, their enforcement points, and the evidence that exercises
them under pressure.}
\label{tab:invariants}
\begin{center}
\small
\begin{tabular}{p{0.30\linewidth}p{0.34\linewidth}p{0.26\linewidth}}
\toprule
\textbf{Invariant} & \textbf{Enforcement} & \textbf{Exercised by} \\
\midrule
1. Never silently wrong --- no path reports success when a verification failed &
every replayed step assertion-gated; outcome gated on authoritative evidence &
0/90 silent failures under drift, decoys, and ambiguity (\S\ref{sec:results-safety}) \\
\addlinespace
2. Never worse than baseline --- plain-agent fallback always reachable and clean &
typed failure classification on every gate; fallback logs \emph{why} it fired &
failed replays fall back and complete cold (\S\ref{sec:results-safety}) \\
\addlinespace
3. Never cross environments &
structural fingerprint mismatch short-circuits before any fuzzy matching &
matcher and continuation cross-environment cases fail with zero dispatches \\
\addlinespace
4. Everything versioned --- store mutations append-only, no in-place edits &
event log, playbook library, site memory, and checkpoints are append-only with rollback &
interrupted writes recover without in-place edits (\S\ref{sec:results-learning}) \\
\addlinespace
5. Every model call tagged &
single client wrapper with a closed source vocabulary; untagged calls fail lint &
all token accounting in \S\ref{sec:results} decomposes exactly by source \\
\bottomrule
\end{tabular}
\end{center}
\end{table}

Two consequences are worth stating. First, invariants 1--3 mean Rote's failure mode is
economic, not behavioral: when memory cannot be trusted, the system pays the cold price
it would have paid anyway, plus a cheap gate check. Second, invariants 4--5 make every
claim in this paper auditable --- each reported number decomposes into tagged calls over
append-only records that can be replayed.

\section{Experimental setup}
\label{sec:setup}
% Benchmark B1--B6 (docs/03), exact external oracles, frozen fixtures + digest-pinned WordPress page
% (T2: 22,279 tokens, zero range over 15 sessions). Baselines: Browser Use 0.13.7 (qualified, T24);
% Stagehand/Skyvern/Magnitude stopped at qualification (T22/T23/T27) -- report why, no ranking.
% Metrics: logical input tokens, billed cost, latency, exact success, silent-failure rate.
\textbf{[TODO]}

\section{Results}
\label{sec:results}
\subsection{Warm-run economics}
% T25: 18/18 exact per harness; tokens -83.1% [82.1, 83.9], cost -68.2% [66.3, 69.9],
% latency -43.6% [39.0, 48.1]. T10: 37.2% slower logical-input growth on the curve.
\subsection{Drift and safety}
\label{sec:results-safety}
% T21: 72/72 repaired exact successes, 0/90 silent failures, 18/18 ambiguity fallbacks, zero repair
% tokens. T35: contract drift stops with BROWSER_CONTRACT_MISMATCH and zero dispatches while cosmetic,
% rename, and wrapper drift keep replaying. T40: B6 lookalike never selected; forced replay dispatches
% nothing.
\subsection{Learning across runs}
\label{sec:results-learning}
% T36: record -> distill -> replay 0 LLM calls on B1/B2. T37: continuation across 2 process restarts,
% no repeated dispatch. [TODO billed: T0 >=80% reduction at parity with automated distillation;
% T2 >=30% generalization from the site brief; >=50% warm steps off the frontier at parity.]
\subsection{Predictor}
\label{sec:results-predictor}
% T38/T39: 99.4% kind+target over 1,520 warm steps (Wilson 98.9-99.7); >=0.9 confidence covers 62.4%
% at 99.8% precision; under-calibrated on fixtures (0.33 -> 99%) -> shadow only. Fig: reliability
% diagram + coverage curve from docs/testing/data/T39-predictor-simulation.json.
\subsection{Long sessions}
% T34: 60-transition SPA, visible history <=31 actions, 3 compaction boundaries, 15/15 pass.
\textbf{[TODO]}

\section{Limitations}
% Synthetic enterprise fixtures for most safety results; one real page family (WordPress); baselines
% partly stopped at qualification; predictor calibration only on fixtures; billed campaign scope.
\textbf{[TODO]}

\section{Conclusion}
\textbf{[TODO]}

\subsubsection*{Reproducibility statement}
% Point to anonymized code + fixtures + T-reports (docs/testing/T*.md, raw data under
% docs/testing/data/), one-command reproduction (T17), pinned baseline versions.
\textbf{[TODO]}

\subsubsection*{AI use statement}
% Required by ICLR 2027 policy. Describe use of LLM coding assistants in implementation/writing.
\textbf{[TODO]}

\bibliography{rote}
\bibliographystyle{iclr2027_conference}

\appendix
\section{Benchmark tasks B1--B6 and oracles}
\section{Action-contract schema and drift taxonomy}
\section{Full test log index}
% Table mapping T1..T40 -> claim used in the paper.

\end{document}
